\chapter*{前言}
\addcontentsline{toc}{chapter}{前言}
\markboth{前言}{前言}

这是一本给初学者写的 Linux 书。

它假设你没有写过程序，没有用过命令行，也没有装过操作系统。它不假设你笨，也不打算把每一件事都讲成三步操作。命令行值得理解，而不是背诵；系统值得理解，而不是靠重装解决问题。

\section*{为什么第一台 Linux 用 NixOS}

多数入门书会选 Ubuntu 或 Debian，这是稳妥的选择：教程多，遇到问题容易搜到答案。本书选 NixOS，理由不是它更流行，而是它把 Linux 系统的构成摆在明面上。

在传统发行版里，系统是几百次安装、修改、升级累积出来的结果。半年之后，没人能完整说出机器上装了什么、改过哪个配置文件。NixOS 把整个系统写成一份配置：装了哪些软件，开了哪些服务，用户是谁，网络怎么配。改配置，重建，切换；出错就回滚到上一代。这件事对初学者格外重要——你可以放心地把系统改坏，因为改坏是可以撤销的。

代价也要说清楚：NixOS 的学习曲线比 Ubuntu 陡，网上大多数教程不能照抄，部分软件需要自己打包。本书会在每一步解释“为什么这样做”，而不是只给一串命令。

\section*{这本书不做什么}

它不试图覆盖 Linux 的全部。内核调优、集群编排、安全审计、嵌入式开发，都不在范围内。它也不承诺让你在读完以后成为专家；它承诺的是：你能独立装好一台 NixOS，能读懂并修改系统配置，能在服务器上部署和排查一个服务，并且知道遇到问题时该往哪里找答案。

\section*{关于 AI}

本书最后一部分讲 AI Agent。它不讨论“AI 会不会取代程序员”，只回答更具体的问题：怎么在自己的机器上装一个 Agent，怎么给它权限，怎么控制成本，怎么审查它改过的文件。工具会变，这些判断不会。

\vspace{12pt}
\begin{flushright}
\bookauthor\\
\bookdate
\end{flushright}
